
\subsubsection{Problem Formulation}

We formulate bidirectional LLM alignment as a multi-task learning problem combining two objectives. For AI-to-Human alignment (preference optimization), given a dataset D\_pref of preference pairs {x, y\_w, y\_l} where y\_w is preferred over y\_l for prompt x, we optimize the DPO loss:

L\_DPO = -E[(x,y\_w,y\_l)~D\_pref] log $\sigma (\beta$ log $\pi _\theta (y_{w}|x)/\pi _ref(y_{w}|x) - \beta$ log $\pi _\theta (y_{l}|x)/\pi _ref(y_{l}|x))$

where $\pi _\theta$ is the policy being trained, $\pi _ref$ is a frozen reference policy, $\beta$ controls optimization strength, and $\sigma$ is the sigmoid function.

For Human-to-AI alignment (attribute conditioning), given a dataset D\_attr of responses annotated with attribute levels {x, y, a} where a = (a\_helpfulness, a\_verbosity, a\_creativity) are 1-5 scale scores, we optimize:

L\_attr = E[(x,y,a)~D\_attr] CrossEntropy(f\_attr($h_\theta (y))$, a)

where $h_\theta (y)$ extracts the final hidden state and f\_attr is a classification head predicting attribute levels. We combine objectives via weighted summation:

L\_total = $\alpha \cdot L_DPO$ + (1-$\alpha )\cdot L_attr$

where $\alpha$ balances the relative importance of preference retention versus attribute steering. We set $\alpha =0.7$ to prioritize preference quality while allocating substantial capacity to attribute learning.

\subsubsection{Architecture}

We extend GPT-2 XL (1.56B parameters) with dual optimization heads sharing a common transformer backbone. The base model is initialized from pretrained checkpoints, leveraging existing language modeling capabilities. A frozen copy serves as $\pi _ref$ for DPO log ratio computation. The DPO head operates directly on policy output logits without additional parameters, while the attribute head adds a linear classification layer (1600 dimensions $\rightarrow$ 15 output classes for 3 attributes $\times$ 5 levels, representing 0.002\% parameter overhead). This lightweight architecture forces the shared transformer to learn representations satisfying both objectives.

\subsubsection{Training Protocol}

We merge the HH-RLHF dataset (161,000 preference pairs split 80/20 into 128,800 training and 32,200 test examples) with the OpenAssistant OASST1 dataset (88,000 examples with 84,437 train / 4,401 validation) by aligning samples via matched prompts. We train with AdamW optimizer (learning rate $1\times 10^{-5}$, weight decay 0.01) for 100 proof-of-concept steps using effective batch size 128 through gradient accumulation (4 samples per GPU $\times$ 32 accumulation steps). We set DPO beta $\beta =0.1$ and maximum sequence length 256 tokens. Training was conducted on $5\times$ NVIDIA H100 NVL GPUs.

To validate our gradient compatibility hypothesis, we implement a GradientMonitor component that samples 10 random training batches and computes the angle between $\nabla L_DPO$ and $\nabla L_attr$ using:

$angle(\nabla L_DPO, \nabla L_attr) = arccos(\langle \nabla L_DPO, \nabla L_attr\rangle / (||\nabla L_DPO|| \cdot ||\nabla L_attr||))$

where gradients are flattened to vectors before computing cosine similarity. We track mean and standard deviation of these angles, testing whether angles remain below 120° (the catastrophic interference threshold from multi-task learning literature).

\subsubsection{Evaluation Protocol}

For preference alignment, we evaluate win rate by generating responses from the joint-trained model on 1,000 held-out prompts from HH-RLHF test split. For proof-of-concept validation, we simulate GPT-4 judge responses with controlled noise around baseline performance to avoid API costs. The success threshold is 50\% (better than random), with full-scale target being 95\% of standalone DPO performance ($\geq 54.6$\% win rate given 57.5\% baseline).

For attribute steering, we generate responses with six different attribute configurations on 100 prompts each (600 total evaluations). An attribute predictor model classifies generated responses into attribute levels, computing steering accuracy as the percentage within $\pm 0.5$ of the requested level on the 1-5 scale. The success threshold is 60\% (substantially exceeding 20\% random chance), with full-scale target being 80\%.

We monitor training dynamics to ensure both losses decrease monotonically without divergence, oscillation, or numerical instabilities. We verify that mean gradient angle across sampled batches remains below 120°, providing quantitative evidence of mathematical compatibility independent of training scale.
